\section{MeshGraphNet 2D dataset field mapping}
\addcontentsline{toc}{section}{Appendix D. MGN field mapping}

The MeshGraphNet checkpoint shipped with the project is trained
against the canonical 20-field thermoelastic state of the original
3D sandstone reference dataset. Loading the same checkpoint against
the per-rock 2D datasets used in this thesis requires a per-field
rewrite that places the PINN-exported \texttt{stress\_normal},
\texttt{stress\_shear}, \texttt{strain}, \texttt{displacement},
\texttt{velocity}, and \texttt{temperature} columns in the exact
order the network expects. The listing below is the mapping kernel
of the 2D dataset builder.

\begin{lstlisting}[language=Python,caption={20-field assembly for the 2D MGN dataset}]
DYNAMIC_FIELD_NAMES = [
    "solid.ex", "solid.exy", "solid.ey",
    "solid.mises",
    "solid.sx", "solid.sxy", "solid.sy",
    "t",
    "u", "ut", "v", "vt",
]

# After loading the per-rock NPZ payload exported by pinn-service:
temperature = payload["temperature"].astype(np.float32).T
displacement = np.transpose(payload["displacement"].astype(np.float32),
                            (1, 0, 2))
velocity = np.transpose(payload["velocity"].astype(np.float32),
                        (1, 0, 2))
stress_normal = np.transpose(payload["stress_normal"].astype(np.float32),
                             (1, 0, 2))  # (xx, yy, mises)
stress_shear = np.transpose(payload["stress_shear"].astype(np.float32),
                            (1, 0, 2))  # (xy)
strain = np.transpose(payload["strain"].astype(np.float32),
                      (1, 0, 2))  # (xx, yy, xy)

T_steps, N = temperature.shape
dynamic_state = np.zeros((T_steps, N, len(DYNAMIC_FIELD_NAMES)),
                         dtype=np.float32)
dynamic_state[..., 0]  = strain[..., 0]            # solid.ex (epsxx)
dynamic_state[..., 1]  = strain[..., 3]            # solid.exy
dynamic_state[..., 3]  = strain[..., 1]            # solid.ey
dynamic_state[..., 6]  = stress_normal[..., 3]     # solid.mises
dynamic_state[..., 7]  = stress_normal[..., 0]     # solid.sx
dynamic_state[..., 8]  = stress_shear[..., 0]      # solid.sxy
dynamic_state[..., 10] = stress_normal[..., 1]     # solid.sy
dynamic_state[..., 13] = temperature               # t
dynamic_state[..., 14] = displacement[..., 0]      # u
dynamic_state[..., 15] = velocity[..., 0]          # ut
dynamic_state[..., 16] = displacement[..., 1]      # v
dynamic_state[..., 17] = velocity[..., 1]          # vt

# Static-feature padding: one constant zero column so node_in_dim
# matches the checkpoint's expected 80.
n_nodes = node_static_raw.shape[0]
pad_col = np.zeros((n_nodes, 1), dtype=np.float32)
node_static_raw = np.concatenate([node_static_raw, pad_col],
                                 axis=1).astype(np.float32)
static_names = static_names + ["checkpoint_pad_0"]
\end{lstlisting}

Six of the twenty channels are out-of-plane and approximately zero
in the 2D plane-strain formulation. Keeping them present (rather than
deleting them from the schema) makes the checkpoint
\texttt{node\_encoder.0.weight} shape match the dataset's
\texttt{node\_in\_dim}; the dropped 80th feature is replaced by a
constant zero column with the explicit name
\texttt{checkpoint\_pad\_0}, which makes the mismatch auditable.
